% !TeX root = ../main.tex
% ---------------------------------------------------------------------------
% Chapter imported from Weekly Study Notes/ Chapter~\ref{ch:rock-mass-parameters}.tex
% Initial enrichment pass: copied directly from the weekly study-note body.
% ---------------------------------------------------------------------------
% ======================================================================
\chapter{Rock Mass Failure Criteria and Design Parameters}
\label{ch:rock-mass-parameters}
% ======================================================================

\section{The generalised Hoek--Brown criterion}
\label{sec:the-generalised-hoek-brown-criterion}

\subsection{Generalised Hoek--Brown Criterion}
% ======================================================================

\paragraph{Rock Mass Failure Criterion}

The generalised Hoek--Brown criterion for jointed rock masses is:
\[
\sigma_1^f =
\sigma_3 +
\sigma_{ci}
\left(
m_b\frac{\sigma_3}{\sigma_{ci}} + s
\right)^a
\]

where:
\begin{itemize}
 \item \(\sigma_1^f\) = major principal stress at failure;
 \item \(\sigma_3\) = confining stress used in the failure criterion;
 \item \(\sigma_{ci}\) = uniaxial compressive strength of intact rock;
 \item \(m_b\), \(s\), \(a\) = rock mass Hoek--Brown parameters;
 \item \(m_i\) = intact rock material constant;
 \item \(GSI\) = Geological Strength Index;
 \item \(D\) = disturbance factor.
\end{itemize}

\begin{keybox}[From intact rock to rock mass]
The intact Hoek--Brown criterion uses \(m_i\), \(s=1\), and \(a=0.5\). 
The generalised Hoek--Brown criterion modifies these values using \(GSI\) and \(D\)
to represent a jointed and disturbed rock mass.
\end{keybox}

% ----------------------------------------------------------------------
\paragraph{Hoek--Brown Rock Mass Parameters}

The rock mass constants are calculated as:
\[
m_b =
m_i
\exp
\left(
\frac{GSI-100}{28-14D}
\right)
\]

\[
s =
\exp
\left(
\frac{GSI-100}{9-3D}
\right)
\]

\[
a =
\frac{1}{2}
+
\frac{1}{6}
\left(
e^{-GSI/15}
-
e^{-20/3}
\right)
\]

\begin{calcbox}[Key check]
For intact rock:
\[
GSI=100,\qquad D=0
\]
then:
\[
m_b=m_i,\qquad s=1,\qquad a=0.5
\]
and the criterion reduces to the original intact Hoek--Brown criterion.
\end{calcbox}

% ======================================================================
\subsection{Disturbance Factor \(D\)}
% ======================================================================

The disturbance factor accounts for blast damage, stress relaxation, and excavation-induced
disturbance around an opening. It reduces the rock mass strength and deformation modulus.

\begin{table}[H]
\centering
\small
\begin{tabularx}{\textwidth}{l Y c}
\toprule
\textbf{Excavation method} & \textbf{Description} & \textbf{Suggested \(D\)} \\
\midrule
TBM excavation & Mechanical excavation, minimal disturbance to surrounding rock & 0 \\
Good quality drill-and-blast & Controlled blasting, smooth wall, presplit & 0.5 \\
Poor quality drill-and-blast & Uncontrolled production blasting, significant damage & 0.8 \\
Very poor blasting & Heavy damage, stress relaxation & 1.0 \\
\bottomrule
\end{tabularx}
\caption{Typical disturbance factor \(D\) for different excavation methods.}
\label{tab:Dfactor}
\end{table}

\begin{warnbox}[Engineering effect of D]
Increasing \(D\) decreases \(m_b\), decreases \(s\), reduces rock mass strength, and reduces
rock mass modulus. Poor blasting can therefore significantly increase support demand even
when the original geological conditions are the same.
\end{warnbox}

% ======================================================================
\section{Calculation workflow and notation}
\label{sec:calculation-workflow-and-notation}

\subsection{ Chapter~\ref{ch:rock-mass-parameters} Calculation Workflow}
% ======================================================================

\begin{figure}[H]
\centering
\resizebox{\textwidth}{!}{%
\begin{tikzpicture}[node distance=0.35cm, font=\small]

\node[smallworkflowbox=orange] (input)
{\textbf{1. Inputs}\\
\footnotesize \(\sigma_{ci},E_i,m_i,GSI,D\)};

\node[smallworkflowbox=green, right=of input] (hb)
{\textbf{2. HB parameters}\\
\footnotesize \(m_b,s,a\)};

\node[smallworkflowbox=blue, right=of hb] (strength)
{\textbf{3. Strength}\\
\footnotesize \(\sigma_{cm},\sigma_{1m}^f,\sigma_{tm}\)};

\node[smallworkflowbox=purple, right=of strength] (modulus)
{\textbf{4. Modulus}\\
\footnotesize \(E_m\)};

\node[smallworkflowbox=teal, right=of modulus] (mc)
{\textbf{5. Equivalent MC}\\
\footnotesize \(c,\phi\)};

\draw[workflowarrow] (input) -- (hb);
\draw[workflowarrow] (hb) -- (strength);
\draw[workflowarrow] (strength) -- (modulus);
\draw[workflowarrow] (modulus) -- (mc);

\end{tikzpicture}
}
\caption{Recommended Chapter~\ref{ch:rock-mass-parameters} rock mass parameter derivation workflow.}
\label{fig:week4_workflow}
\end{figure}

\begin{notebox}[Spreadsheet advice]
Use one row for each rock type and one separate block for each disturbance scenario
(\(D=0\) and \(D=0.8\)). All equations should reference input cells rather than pasted
values.
\end{notebox}

% ======================================================================
\subsection{Stress Notation Used Chapter~\ref{ch:rock-mass-parameters}}
% ====================================================================== Chapter~\ref{ch:rock-mass-parameters} uses stress symbols in two related but different ways. This distinction follows
 Chapter~\ref{ch:stress-intact-rock}.

\begin{table}[H]
\centering
\small
\begin{tabularx}{\textwidth}{l l Y}
\toprule
\textbf{Symbol} & \textbf{Meaning} & \textbf{Use Chapter~\ref{ch:rock-mass-parameters}} \\
\midrule
\(\sigma_v\) & Vertical in-situ stress & Calculated from \(\gamma z\). \\
\(\sigma_{h,K}\) & Horizontal stress from project \(K\) ratio & Calculated from \(K\sigma_v\). \\
\(\sigma_1^{IS}\) & Maximum in-situ principal stress & Current maximum stress demand. \\
\(\sigma_3^{IS}\) & Minimum in-situ principal stress & Used as the in-situ confinement for \(\sigma_{1m}^f\). \\
\(\sigma_3\) in HB & Variable confining stress in the failure criterion & Used to generate the HB envelope. \\
\(\sigma_1^f\) & Major principal stress at failure & Strength capacity predicted by HB. \\
\(\sigma_{3,\mathrm{fit,max}}\) & Upper fitting limit for HB--MC conversion & Used only to define the fitting range. \\
\bottomrule
\end{tabularx}
\caption{Stress notation used Chapter~\ref{ch:rock-mass-parameters}.}
\label{tab:stress_notation}
\end{table}

\begin{keybox}[Most important distinction]
\[
\sigma_3^{IS}=\min(\sigma_v,\sigma_{h,K})
\]
is the simplified minimum in-situ principal stress.

\[
\sigma_{3,\mathrm{fit,max}}=\max(\sigma_v,\sigma_{h,K})
\]
is the upper fitting limit used to convert the nonlinear Hoek--Brown envelope to
equivalent Mohr--Coulomb parameters.

They are not the same quantity.
\end{keybox}

% ======================================================================
\section{Rock mass strength and deformability}
\label{sec:rock-mass-strength-and-deformability}

\subsection{Rock Mass Strength Parameters}
% ======================================================================

Once \(m_b\), \(s\), and \(a\) are known, the following strength parameters can be derived.

% ----------------------------------------------------------------------
\paragraph{Uniaxial Compressive Strength of Rock Mass}

Set \(\sigma_3=0\) in the generalised Hoek--Brown equation:
\[
\sigma_{cm}
=
\sigma_{ci}s^a
\]

where \(\sigma_{cm}\) is the uniaxial compressive strength of the rock mass.

\begin{notebox}[Interpretation]
\(\sigma_{cm}\) is usually much smaller than \(\sigma_{ci}\), especially for low GSI or high
disturbance factor \(D\). This reflects the weakening effect of discontinuities and
excavation disturbance.
\end{notebox}

% ----------------------------------------------------------------------
\paragraph{Triaxial Strength at In-Situ Confinement}

For the design study:
\[
\sigma_v = \gamma z
\]
\[
\sigma_{h,K}=K\sigma_v
\]

The simplified in-situ minor principal stress is:
\[
\sigma_3^{IS}=\min(\sigma_v,\sigma_{h,K})
\]

The triaxial rock mass strength at this in-situ confinement is:
\[
\sigma_{1m}^f =
\sigma_3^{IS}
+
\sigma_{ci}
\left(
m_b\frac{\sigma_3^{IS}}{\sigma_{ci}} + s
\right)^a
\]

\begin{warnbox}[Do not confuse demand and capacity]
\(\sigma_3^{IS}\) is used as the confining stress level. 
\(\sigma_{1m}^f\) is the major principal stress required to cause failure at that confinement.
It is a strength capacity, not the current in-situ maximum stress.
\end{warnbox}

% ----------------------------------------------------------------------
\paragraph{Tensile Strength of Rock Mass}

From the Hoek--Brown criterion:
\[
\sigma_{tm}
=
-\frac{s\sigma_{ci}}{m_b}
\]

The negative sign indicates tensile stress under the compression-positive convention.

An alternative estimate using the Hoek--Martin tensile cut-off is:
\[
\sigma_{tm,HM}
=
-\frac{\sigma_{cm}}{8.62+0.7m_i}
\]

\begin{notebox}[Why compare two tensile estimates?]
The direct Hoek--Brown tensile strength and the Hoek--Martin cut-off may differ.
For design discussion, compare both values and explain which one appears more conservative
or more realistic for brittle cracking.
\end{notebox}

% ======================================================================
\subsection{Deformation Modulus of Rock Mass}
% ======================================================================

The rock mass deformation modulus \(E_m\) is a key input for deformation analysis,
support design, and numerical modelling. Unlike intact Young's modulus \(E_i\), the
rock mass modulus accounts for discontinuities, weathering, blockiness, joint condition,
and excavation disturbance.

\begin{keybox}[Important distinction]
\[
E_i = \text{intact rock Young's modulus from laboratory testing}
\]
\[
E_m = \text{equivalent deformation modulus of the jointed rock mass}
\]
Normally:
\[
E_m < E_i
\]
especially when the rock mass has low GSI, persistent joints, weathering, or excavation
disturbance.
\end{keybox}

Several empirical equations exist for estimating \(E_m\). They may give different results
because they were developed from different datasets and classification systems. Therefore,
they should be treated as estimates rather than exact material constants.

% ----------------------------------------------------------------------
\paragraph{Hoek--Diederichs Generalised Equation}

When laboratory intact modulus \(E_i\) is available, the generalised Hoek--Diederichs
equation is preferred:

\[
E_m =
E_i
\left(
0.02+
\frac{1-D/2}
{1+e^{(60+15D-GSI)/11}}
\right)
\]

where:
\begin{itemize}
 \item \(E_m\) = rock mass deformation modulus;
 \item \(E_i\) = intact rock Young's modulus;
 \item \(GSI\) = Geological Strength Index;
 \item \(D\) = disturbance factor.
\end{itemize}

\begin{notebox}[Units]
The output unit follows the unit used for \(E_i\). If \(E_i\) is entered in GPa, then
\(E_m\) is obtained in GPa. If \(E_i\) is entered in MPa, then \(E_m\) is obtained in MPa.
\end{notebox}

\begin{keybox}[Preferred for design]
Because the design study provides \(E_i\) for each rock type, this generalised
Hoek--Diederichs equation should be used as the \textbf{primary} estimate of \(E_m\).
\end{keybox}

% ----------------------------------------------------------------------
\paragraph{Hoek--Diederichs Simplified Equation}

When reliable intact modulus \(E_i\) is not available, the simplified Hoek--Diederichs
equation may be used:

\[
E_m~(\text{MPa}) =
100000
\left(
\frac{1-D/2}
{1+e^{(75+25D-GSI)/11}}
\right)
\]

or, equivalently, in GPa:

\[
E_m~(\text{GPa}) =
100
\left(
\frac{1-D/2}
{1+e^{(75+25D-GSI)/11}}
\right)
\]

\begin{warnbox}[Correction from common mistake]
The simplified Hoek--Diederichs equation uses:
\[
75+25D-GSI
\]
The generalised Hoek--Diederichs equation uses:
\[
60+15D-GSI
\]
Do not mix these two expressions.
\end{warnbox}

\begin{warnbox}[Limitation]
The simplified equation does not use the measured \(E_i\). It can therefore overestimate
\(E_m\) for weak rocks with low intact modulus, such as mudstone.
\end{warnbox}

% ----------------------------------------------------------------------
\paragraph{Hoek et al. GSI--UCS Equation}

An earlier Hoek et al. equation estimates rock mass modulus using \(GSI\), disturbance
factor \(D\), and intact compressive strength \(\sigma_{ci}\):

\[
E_m~(\text{GPa}) =
\left(1-\frac{D}{2}\right)
\sqrt{\frac{\sigma_{ci}}{100}}
\,
10^{(GSI-10)/40}
\qquad
(\sigma_{ci} \leq 100~\text{MPa})
\]

For stronger intact rocks, it is commonly capped as:

\[
E_m~(\text{GPa}) =
\left(1-\frac{D}{2}\right)
10^{(GSI-10)/40}
\qquad
(\sigma_{ci} > 100~\text{MPa})
\]

\begin{notebox}[Use as a comparison only]
This equation is useful when \(E_i\) is not available but \(\sigma_{ci}\) and GSI are known.
For the design study, it should be used as a comparison only because measured \(E_i\) values
are provided.
\end{notebox}

% ----------------------------------------------------------------------
\paragraph{RMR-Based Equations}

Some older empirical equations estimate \(E_m\) directly from RMR.

\subsubsection{Bieniawski Equation}

For better-quality rock masses, a simple linear estimate is:

\[
E_m~(\text{GPa}) = 2RMR - 100
\qquad
(RMR > 50)
\]

\subsubsection{Serafim--Pereira Equation}

Another widely used RMR-based estimate is:

\[
E_m~(\text{GPa}) =
10^{(RMR-10)/40}
\]

\begin{warnbox}[RMR-based equations]
RMR-based equations are simple, but they do not explicitly include intact modulus
\(E_i\), disturbance factor \(D\), or excavation method. They should be used only as
rough comparison checks.
\end{warnbox}

% ----------------------------------------------------------------------
\paragraph{Q-Based Equation}

A Q-system-based estimate can be written using the normalised \(Q_c\) value:

\[
E_m~(\text{GPa}) =
10Q_c^{1/3}
\]

where:

\[
Q_c =
Q\frac{\sigma_{ci}}{100}
\]

\begin{notebox}[When this may be useful]
The Q-based modulus estimate may be useful when Q-system data are more reliable than
RMR or GSI, especially in tunnel support design. However, it should still be compared
with GSI-based estimates.
\end{notebox}

% ----------------------------------------------------------------------
\paragraph{Recommended Approach for design}

Use the following hierarchy:

\begin{enumerate}
 \item \textbf{Primary estimate:} Hoek--Diederichs generalised equation, because \(E_i\),
 GSI, and \(D\) are available.
 
 \item \textbf{Secondary comparison:} Hoek--Diederichs simplified equation.
 
 \item \textbf{Optional sensitivity checks:} Hoek et al. GSI--UCS equation,
 RMR-based equations, or Q-based equation.
\end{enumerate}

\begin{keybox}[Recommended reporting wording]
``The generalised Hoek--Diederichs equation is adopted as the primary estimate because it
uses the measured intact modulus \(E_i\), the adopted GSI, and the disturbance factor
\(D\). Other empirical equations are used only as sensitivity checks because they do not
fully account for the project-specific intact modulus or excavation disturbance.''
\end{keybox}

% ----------------------------------------------------------------------
\paragraph{Suggested Comparison Table}

\begin{table}[H]
\centering
\small
\begin{tabularx}{\textwidth}{l c c c c}
\toprule
\textbf{Method} & \textbf{Granite} & \textbf{Sandstone} & \textbf{Mudstone} & \textbf{Limestone} \\
\midrule
Hoek--Diederichs generalised & & & & \\
Hoek--Diederichs simplified & & & & \\
Hoek et al. GSI--UCS equation & & & & \\
Bieniawski RMR equation & & & & \\
Serafim--Pereira RMR equation & & & & \\
Q-based equation & & & & \\
\midrule
\textbf{Adopted \(E_m\)} & & & & \\
\bottomrule
\end{tabularx}
\caption{Suggested comparison of rock mass deformation modulus estimates. Complete separately for \(D=0\) and \(D=0.8\) where applicable.}
\label{tab:modulus_comparison}
\end{table}
% ======================================================================
\section{Converting Hoek--Brown to Equivalent Mohr--Coulomb Parameters}
% ======================================================================

Many design calculations and numerical models require Mohr--Coulomb parameters
\(c\) and \(\phi\). Because Hoek--Brown is nonlinear, \(c\) and \(\phi\) are obtained by
fitting an equivalent linear envelope over a selected stress range.

\paragraph{Selecting the Fitting Range}

Use:
\[
0 \leq \sigma_3 \leq \sigma_{3,\mathrm{fit,max}}
\]

For the design study:
\[
\sigma_{3,\mathrm{fit,max}}=\max(\sigma_v,\sigma_{h,K})
\]

\begin{warnbox}[Stress range matters]
The equivalent Mohr--Coulomb parameters are not unique. A smaller fitting range usually
gives a higher \(\phi\) and lower \(c\), while a larger fitting range usually gives a lower
\(\phi\) and higher \(c\). Always state the fitting range used.
\end{warnbox}

\begin{notebox}[Why use sigma3 fit max?]
\(\sigma_{3,\mathrm{fit,max}}\) is the upper confining stress used to fit the HB envelope.
It is not the minimum in-situ principal stress. It is a modelling choice for obtaining
equivalent \(c\) and \(\phi\).
\end{notebox}

% ----------------------------------------------------------------------
\paragraph{Method 1: Mohr Circle Regression}

\begin{enumerate}
 \item Select a series of \(\sigma_3\) values between 0 and \(\sigma_{3,\mathrm{fit,max}}\).
 \item Compute corresponding \(\sigma_1^f\) values using:
 \[
 \sigma_1^f =
 \sigma_3+
 \sigma_{ci}
 \left(
 m_b\frac{\sigma_3}{\sigma_{ci}}+s
 \right)^a
 \]
 \item For each pair, compute:
 \[
 C=\frac{\sigma_1^f+\sigma_3}{2}
 \]
 \[
 R=\frac{\sigma_1^f-\sigma_3}{2}
 \]
 \item Fit:
 \[
 R=A+BC
 \]
 \item Convert:
 \[
 \phi=\arcsin(B)
 \]
 \[
 c=\frac{A}{\cos\phi}
 \]
\end{enumerate}

\begin{calcbox}[Excel functions]
\[
B=\texttt{SLOPE(R\_values,C\_values)}
\]
\[
A=\texttt{INTERCEPT(R\_values,C\_values)}
\]
\[
R^2=\texttt{RSQ(R\_values,C\_values)}
\]
\end{calcbox}

% ----------------------------------------------------------------------
\paragraph{Method 2: Linear Regression in \(\sigma_1^f-\sigma_3\) Space}

Using the same generated data, fit:
\[
\sigma_1^f=a_{\mathrm{fit}}+b_{\mathrm{fit}}\sigma_3
\]

Then:
\[
\sin\phi=
\frac{b_{\mathrm{fit}}-1}
{b_{\mathrm{fit}}+1}
\]

\[
c=
\frac{a_{\mathrm{fit}}(1-\sin\phi)}
{2\cos\phi}
\]

\begin{notebox}[Cross-check]
Methods 1 and 2 should give similar \(c\) and \(\phi\) values if the same stress range
and data points are used. Large differences usually indicate a spreadsheet error.
\end{notebox}

% ======================================================================
\section{Worked Example: Sandstone Rock Mass}
\label{sec:ch04-worked-example}
% ======================================================================

Consider an illustrative sandstone rock mass:
\[
GSI=75,\qquad m_i=17,\qquad \sigma_{ci}=75~\text{MPa},\qquad D=0
\]

The calculated Hoek--Brown parameters are approximately:
\[
m_b=6.96,\qquad s=0.062,\qquad a=0.501
\]

Generated values are shown in Table~\ref{tab:example}.

\begin{table}[H]
\centering
\small
\begin{tabular}{c c c c}
\toprule
\(\sigma_3\) (MPa) & \(\sigma_1^f\) (MPa) & \(C\) (MPa) & \(R\) (MPa) \\
\midrule
0 & 18.7 & 9.3 & 9.3 \\
5 & 60.5 & 32.8 & 27.8 \\
10 & 93.9 & 52.0 & 42.0 \\
15 & 124.0 & 69.5 & 54.5 \\
20 & 152.0 & 86.0 & 66.0 \\
\bottomrule
\end{tabular}
\caption{Generated Hoek--Brown data for an illustrative sandstone rock mass.}
\label{tab:example}
\end{table}

Using linear regression in \(\sigma_1^f-\sigma_3\) space:
\[
b_{\mathrm{fit}}\approx 6.67,
\qquad
a_{\mathrm{fit}}\approx 18.7
\]

\[
\sin\phi=
\frac{6.67-1}{6.67+1}
=
0.739
\]

\[
\phi\approx47.7^\circ
\]

\[
c=
\frac{18.7(1-0.739)}
{2\cos(47.7^\circ)}
\approx
3.6~\text{MPa}
\]

\begin{keybox}[Interpretation]
The equivalent Mohr--Coulomb parameters represent a linear approximation to the nonlinear
Hoek--Brown envelope over the selected stress range. They should not be treated as
unique material constants independent of confinement range.
\end{keybox}

% ======================================================================
\section{Design data and calculation tables}
\label{sec:ch4-design-data}

\subsection{Project Input Data for design}
% ======================================================================

The following table summarises the intact rock properties and project stress conditions.
The GSI values must be taken from engineers' Chapter~\ref{ch:classification} results.

\begin{table}[H]
\centering
\scriptsize
\begin{tabularx}{\textwidth}{l c c c c}
\toprule
\textbf{Parameter} & \textbf{Granite} & \textbf{Sandstone} & \textbf{Mudstone} & \textbf{Limestone} \\
\midrule
\(\sigma_{ci}\) (MPa) & 265 & 110 & 27 & 85 \\
\(E_i\) (GPa) & 80 & 30 & 0.76 & 20 \\
\(m_i\) & 32 & 17 & 4 & 12 \\
Depth \(z\) (m) & 780 & 450 & 100 & 80 \\
\(K=\sigma_{h,K}/\sigma_v\) & 3.2 & 1.8 & 1.0 & 2.0 \\
\(\gamma\) (kN/m\(^3\)) & 27 & 27 & 27 & 27 \\
\midrule
\(\sigma_v\) (MPa) & 21.06 & 12.15 & 2.70 & 2.16 \\
\(\sigma_{h,K}\) (MPa) & 67.39 & 21.87 & 2.70 & 4.32 \\
\(\sigma_1^{IS}=\max(\sigma_v,\sigma_{h,K})\) (MPa) & 67.39 & 21.87 & 2.70 & 4.32 \\
\(\sigma_3^{IS}=\min(\sigma_v,\sigma_{h,K})\) (MPa) & 21.06 & 12.15 & 2.70 & 2.16 \\
\(\sigma_{3,\mathrm{fit,max}}\) (MPa) & 67.39 & 21.87 & 2.70 & 4.32 \\
\midrule
GSI Chapter~\ref{ch:classification} & \underline{\hspace{1.5cm}} & \underline{\hspace{1.5cm}} & \underline{\hspace{1.5cm}} & \underline{\hspace{1.5cm}} \\
\bottomrule
\end{tabularx}
\caption{Input data Chapter~\ref{ch:rock-mass-parameters} rock mass parameter derivation.}
\label{tab:inputdata}
\end{table}

% ======================================================================
\subsection{Recommended Calculation Tables for Design Reports}
% ======================================================================

\paragraph{Hoek--Brown Parameter Template}

\begin{table}[H]
\centering
\small
\begin{tabularx}{0.95\textwidth}{l c c c c}
\toprule
\textbf{Parameter} & \textbf{Granite} & \textbf{Sandstone} & \textbf{Mudstone} & \textbf{Limestone} \\
\midrule
\(GSI\) & & & & \\
\(D\) & & & & \\
\(m_i\) & & & & \\
\(m_b\) & & & & \\
\(s\) & & & & \\
\(a\) & & & & \\
\bottomrule
\end{tabularx}
\caption{Suggested table for Hoek--Brown rock mass parameters. Complete separately for \(D=0\) and \(D=0.8\).}
\label{tab:hb_template}
\end{table}

\paragraph{Strength and Modulus Template}

\begin{table}[H]
\centering
\small
\begin{tabularx}{0.95\textwidth}{l c c c c}
\toprule
\textbf{Parameter} & \textbf{Granite} & \textbf{Sandstone} & \textbf{Mudstone} & \textbf{Limestone} \\
\midrule
\(\sigma_{cm}\) (MPa) & & & & \\
\(\sigma_{1m}^f\) at \(\sigma_3^{IS}\) (MPa) & & & & \\
\(\sigma_{tm}\) from HB (MPa) & & & & \\
\(\sigma_{tm,HM}\) from cut-off (MPa) & & & & \\
\(E_m\), generalised (GPa) & & & & \\
\(E_m\), simplified (GPa) & & & & \\
\bottomrule
\end{tabularx}
\caption{Suggested table for rock mass strength and modulus parameters. Complete separately for \(D=0\) and \(D=0.8\).}
\label{tab:strength_template}
\end{table}

\paragraph{Equivalent Mohr--Coulomb Template}

\begin{table}[H]
\centering
\small
\begin{tabularx}{0.95\textwidth}{l c c c c}
\toprule
\textbf{Parameter} & \textbf{Granite} & \textbf{Sandstone} & \textbf{Mudstone} & \textbf{Limestone} \\
\midrule
\(\sigma_{3,\mathrm{fit,max}}\) (MPa) & & & & \\
\(c\), Method 1 (MPa) & & & & \\
\(\phi\), Method 1 (\(^\circ\)) & & & & \\
\(c\), Method 2 (MPa) & & & & \\
\(\phi\), Method 2 (\(^\circ\)) & & & & \\
Difference/comment & & & & \\
\bottomrule
\end{tabularx}
\caption{Suggested table for equivalent Mohr--Coulomb parameters. Complete separately for \(D=0\) and \(D=0.8\).}
\label{tab:mc_template}
\end{table}

% ======================================================================
